\documentclass[conference]{IEEEtran}
\usepackage[utf8]{inputenc}
\usepackage[spanish]{babel}
\usepackage{cite}
\usepackage{amsmath,amssymb,amsfonts}
\usepackage{graphicx}
\usepackage[kw]{pseudo}

\title{Planificación automática aplicada al Senku}

\author{
\IEEEauthorblockN{Gregorio Morillo Zamora}
\IEEEauthorblockA{
\textit{Dpto. Ciencias de la Computación e Inteligencia Artificial}\\
\textit{Universidad de Sevilla}\\
Sevilla, España\\
XVJ8387 gremorzam@alum.us.es}

\and

\IEEEauthorblockN{Raúl Jimenez Pérez}
\IEEEauthorblockA{
\textit{Dpto. Ciencias de la Computación e Inteligencia Artificial}\\
\textit{Universidad de Sevilla}\\
Sevilla, España\\
LSW694 raujimper@alum.us.es}
}

\begin{document}

\maketitle

\begin{abstract}
En este trabajo se aborda la resolución del juego de solitario Senku mediante técnicas de planificación automática. Se explora el modelado del dominio del juego utilizando PDDL y la implementación de algoritmos de búsqueda para encontrar secuencias de acciones válidas que permitan alcanzar un estado objetivo, definido como la presencia de una única pieza en el tablero.
\end{abstract}

\section{Introducción}
El juego del Senku, también conocido como solitario, constituye un entorno clásico y de gran utilidad en el campo de la Inteligencia Artificial. Al tratarse de un juego de un solo jugador, permite la evaluación rigurosa del desempeño de agentes y algoritmos de búsqueda en un espacio de estados definido.

El objetivo principal de este proyecto es el diseño de un sistema capaz de representar el problema del Senku, permitiendo su resolución automática mediante la aplicación de algoritmos de búsqueda. El trabajo se estructura en torno a varios objetivos específicos:
\begin{itemize}
    \item Expresar formalmente el problema del Senku mediante el lenguaje PDDL.
    \item Implementar un motor de búsqueda en Python capaz de resolver el tablero clásico bidimensional.
    \item Aplicar algoritmos específicos de búsqueda, evaluando su eficiencia y capacidad para resolver diversas variantes del tablero.
    \item Documentar el proceso de diseño, las decisiones tomadas y los resultados experimentales obtenidos bajo una metodología científica.
\end{itemize}

A diferencia de versiones más restrictivas del problema, en este trabajo se considera como estado objetivo cualquier configuración final donde exista únicamente una pieza en el tablero, independientemente de su posición final. No obstante, se realiza una exploración adicional de configuraciones iniciales que permitan alcanzar el objetivo en la posición del hueco original, profundizando así en la complejidad y las propiedades del juego.

La memoria se organiza de la siguiente manera: tras esta introducción, se describe el lenguaje de modelado, se detallan los algoritmos implementados, se presentan los resultados obtenidos y, finalmente, se extraen las conclusiones del estudio.

\section{Metodología}

En esta sección se describen los algoritmos implementados para la resolución del Senku. Cada algoritmo se presenta en un subapartado independiente, detallando su funcionamiento, las variantes consideradas y las decisiones de diseño adoptadas para superar las dificultades encontradas durante su implementación.

\subsection{Búsqueda básica en espacio de estados}

El primer algoritmo implementado consiste en una exploración sistemática del árbol de juego del Senku mediante búsqueda ciega, es decir, sin emplear ninguna función heurística que guíe la exploración. El problema se modela como un grafo dirigido donde cada nodo representa una configuración del tablero (posición de las piezas y huecos) y cada arista dirigida representa un movimiento válido, esto es, el salto de una pieza sobre otra adyacente hacia un hueco libre.

El algoritmo mantiene dos estructuras de datos fundamentales durante su ejecución. La primera es la \textit{frontera}, una colección que almacena los nodos pendientes de explorar y cuyo orden de extracción determina la estrategia de búsqueda. La segunda es el \textit{conjunto de visitados}, que almacena una representación resumida (mediante una función hash) de todos los estados que ya han sido explorados, evitando así que el algoritmo entre en ciclos infinitos al alcanzar un estado previamente visitado.

En cada iteración, el algoritmo extrae un nodo de la frontera, genera todos los movimientos válidos desde el estado representado por dicho nodo y, para cada uno, aplica la transición correspondiente para obtener el estado sucesor. Si el estado sucesor no ha sido visitado con anterioridad, se comprueba si cumple la condición de objetivo (una única pieza en el centro del tablero). En caso afirmativo, el algoritmo reconstruye la secuencia de acciones recorriendo hacia atrás las referencias al nodo padre; en caso negativo, el nodo sucesor se añade a la frontera y el original al conjunto de visitados, y el proceso continúa. Si la frontera se vacía sin haber encontrado una solución, el algoritmo concluye que no existe ningún plan de movimientos que alcance el estado objetivo desde la configuración inicial dada.

\subsubsection{Variantes: BFS y DFS}

El algoritmo admite dos variantes según la política de extracción de la frontera:

\begin{itemize}
    \item \textbf{Búsqueda en anchura (BFS)}: la frontera se gestiona como una cola FIFO (\textit{first in, first out}). Esta variante garantiza que la primera solución encontrada es la de menor número de movimientos posibles. Sin embargo, su consumo de memoria crece exponencialmente con la profundidad del árbol, lo que la hace inviable para el tablero clásico completo en un ordenador convencional, dado el enorme volumen de estados que deben mantenerse simultáneamente en la frontera.
    \item \textbf{Búsqueda en profundidad (DFS)}: la frontera se gestiona como una pila LIFO (\textit{last in, first out}). Esta variante requiere menos memoria que BFS, ya que explora un único camino hasta agotarlo antes de retroceder. No obstante, no garantiza encontrar la solución de profundidad mínima y, sin un límite adecuado, puede perderse en ramas muy profundas del árbol sin alcanzar el objetivo.
\end{itemize}

Para la verificación temprana del correcto funcionamiento de ambas variantes se empleó un tablero reducido de solo tres piezas. Esto permitió depurar la lógica del algoritmo y confirmar que los planes generados eran válidos antes de abordar el tablero clásico completo, cuyo espacio de búsqueda es significativamente mayor.

\subsubsection{Pseudocódigo}

A continuación se presenta el pseudocódigo unificado del algoritmo de búsqueda básica. Presentamos antes algunas aclaraciones para facilitar la comprensión del pseudocódigo que define el algoritmo. La función \texttt{elegir\_nodo} determina la variante: si se trata de BFS, extrae el primer elemento de la frontera (FIFO); si se trata de DFS, extrae el último (LIFO). La función \texttt{generar\_movimientos} obtiene todos los saltos válidos desde el estado actual comprobando, para cada tripleta de casillas adyacentes, que la casilla origen contiene una pieza, la casilla intermedia contiene una pieza y la casilla destino está vacía. La función \texttt{aplicar} ejecuta la transición y produce el nuevo estado. La función \texttt{hash} construye un identificador inmutable del estado considerando únicamente aquellas casillas ocupadas por una pieza, lo que permite la detección eficiente de estados repetidos.

\begin{figure}
  \begin{pseudo}*
    \hd{\fn{busquedaBasica}}(estadoInicial, usarBFS) \\*
    \multicolumn{2}{l}{\textbf{Entrada}: \textnormal{estado inicial del tablero, modo de búsqueda}} \\*
    \multicolumn{2}{l}{\textbf{Salida}: plan de acciones o indicador de fracaso} \\
    frontera $\leftarrow$ [\pr{nodo}(estadoInicial)] \\
    visitados $\leftarrow$ \{ \pr{hash}(estadoInicial) \} \\
    \textbf{mientras} frontera \textnormal{no esté vacía} \textbf{hacer} \\+
        nodo $\leftarrow$ \pr{elegirNodo}(frontera, usarBFS) \\
        acciones $\leftarrow$ \pr{generarMovimientos}(nodo.estado) \\
        \textbf{para cada} \textnormal{acción} \textbf{en} \textnormal{acciones} \textbf{hacer} \\+
            sucesor $\leftarrow$ \pr{aplicar}(\textnormal{acción}, nodo.estado) \\
            \textbf{si} \pr{hash}(sucesor) $\notin$ visitados \textbf{entonces} \\+
                \textbf{si} \pr{esObjetivo}(sucesor) \textbf{entonces} \\+
                    \textbf{devolver} \pr{reconstruirPlan}(sucesor) \\-
                visitados $\leftarrow$ visitados $\cup$ \{ \pr{hash}(sucesor) \} \\
                frontera $\leftarrow$ frontera $\cup$ [\pr{nodo}(sucesor, nodo, \textnormal{acción})] \\-
            \textbf{fin si} \\-
        \textbf{fin para} \\-
    \textbf{fin mientras} \\
    \textbf{devolver} \textnormal{fracaso}
  \end{pseudo}
  \caption{Algoritmo de búsqueda básica en espacio de estados (BFS/DFS)}
  \label{pcd:basica}
\end{figure}

\subsubsection{Dificultades encontradas y decisiones de diseño}

La principal dificultad encontrada durante la implementación de este algoritmo fue el rendimiento, motivado por la enorme dimensión del espacio de estados del Senku clásico. A continuación se detallan los problemas específicos y las soluciones adoptadas.

El primer problema identificado fue la generación de los movimientos aplicables desde un estado. La librería Unified Planning proporciona un simulador que, para cada estado, recalca todas las posibles instanciaciones de la acción \texttt{saltar} combinando todos los objetos del dominio, lo que resulta en un proceso de grounded intensivo. Para el tablero clásico, con 33 casillas, esto supone evaluar más de treinta mil combinaciones de parámetros en cada nodo, con un coste temporal del orden de decenas de segundos por operación, lo cual es inasumible para una búsqueda que debe explorar millones de nodos. La solución consistió en extraer las relaciones de adyacencia del problema una única vez al inicio de la búsqueda y, a partir de entonces, generar los movimientos válidos comprobando directamente sobre cada tripleta de casillas adyacentes si se cumplen las precondiciones de la acción. Esta comprobación se reduce a verificar tres condiciones booleanas en memoria, reduciendo el coste drásticamente y haciendo viable la exploración.

El segundo problema derivaba del cálculo del identificador hash de cada estado para el conjunto de visitados. Inicialmente, el hash se construía recorriendo todos los fluentes definidos en el estado inicial del problema PDDL. Sin embargo, el predicado \texttt{adyacente} genera miles de fluentes estáticos (que nunca cambian durante la partida) y que son innecesarios para distinguir un estado de otro. La solución fue filtrar estos fluentes, considerando únicamente aquellos correspondientes a los predicados \texttt{pieza} y \texttt{vacio}, que son los únicos que varían con cada movimiento. Esto redujo el número de fluentes a evaluar por hash en más de un 99\%, acelerando significativamente cada llamada.

Estas optimizaciones permitieron que el algoritmo completase la exploración a una velocidad del orden de varios cientos de nodos por segundo, haciendo factible la resolución del tablero clásico en un tiempo del orden de horas. No obstante, queda patente que una búsqueda puramente ciega, sin información heurística que guíe la exploración hacia los estados más prometedores, tiene dificultades para escalar a problemas de gran tamaño. Esta limitación motiva la incorporación de técnicas de búsqueda informada que se describen en las siguientes secciones.

\section{Descripción del módulo src/pddl\_reader.py}

\subsection{Propósito}
El módulo \texttt{src/pddl\_reader.py} proporciona utilidades para cargar un dominio y un problema PDDL usando la librería \texttt{unified\_planning}, extraer los fluentes iniciales activos y calcular una función de evaluación simple (denominada "función pagoda") sobre un estado dado. Está pensado como una pieza de integración entre los ficheros PDDL en el directorio \texttt{PDDL/} y los algoritmos de búsqueda implementados en \texttt{src/}.

\subsection{Estructura y API principal}
El módulo define la clase \texttt{SenkuPDDLReader} con la siguiente interfaz relevante:

\begin{itemize}
    \item \textbf{\_\_init\_\_}(domain\_path, problem\_path): almacena las rutas del dominio y del problema.
    \item \textbf{read\_problem()}: utiliza \texttt{PDDLReader.parse\_problem} para parsear y cargar el problema PDDL en memoria, guardándolo en \texttt{self.problem}.
    \item \textbf{cargar\_fluentes()}: devuelve un diccionario con los fluentes del estado inicial cuyo valor es verdadero. Internamente usa \texttt{self.problem.initial\_values}.
    \item \textbf{obtener\_acciones\_aplicables()}: devuelve la colección de acciones del problema (\texttt{self.problem.actions}).
    \item \textbf{funcion\_pagoda(fluentes\_activos, variante)}: calcula una suma entera según una asignación binaria de casillas (tres variantes implementadas). Permite evaluar estados con una heurística muy simple usada para experimentación.
\end{itemize}

Además hay métodos auxiliares internos:
\begin{itemize}
    \item \textbf{\_extraer\_coordenadas(fluente)}: normaliza la representación de un fluente de pieza aceptando tanto objetos fluente devueltos por Unified Planning como cadenas del tipo \texttt{pieza(c31)}. Devuelve el identificador de casilla (por ejemplo \texttt{c31}) o \texttt{None} si no aplica.
\end{itemize}

\subsection{Decisiones de diseño importantes}
\begin{itemize}
    \item \textbf{Integración con Unified Planning}: el módulo delega el parseo completo del PDDL en \texttt{unified\_planning.io.PDDLReader}. Esto simplifica el código y garantiza compatibilidad con la API de UP, a costa de depender de su representación interna de fluentes y acciones.
    \item \textbf{Representación flexible de fluentes}: \texttt{\_extraer\_coordenadas} admite tanto objetos fluente como cadenas, lo que facilita reutilizar la función con distintas partes del proyecto (por ejemplo, al comparar valores guardados en estructuras propias frente a la salida directa de UP).
    \item \textbf{Función pagoda ligera}: se implementaron tres variantes sencillas de la función pagoda como heurísticas experimentales. Las variantes están basadas en patrones sobre índices de casilla (par/impar) y son rápidas de calcular.
    \item \textbf{Simplicidad y legibilidad}: el módulo prioriza claridad sobre microoptimización; operaciones de carga y cálculo se mantienen explícitas para facilitar la inspección durante la experimentación.
\end{itemize}

\subsection{Limitaciones y observaciones}
\begin{itemize}
    \item \textbf{Formato de casilla asumido}: \texttt{funcion\_pagoda} asume identificadores de casilla de longitud 3 que comienzan por \texttt{c} y continúan con dos dígitos únicos (por ejemplo \texttt{c31}). Esto falla para índices con más de una cifra por coordenada y debería refactorizarse si se usan tableros con dimensiones mayores o distintas convenciones de nombrado.
    \item \textbf{Ausencia de manejo de errores robusto}: el módulo captura excepciones en el bloque \texttt{if __name__ == "__main__"} sólo para imprimir errores; sería recomendable añadir validaciones y mensajes más informativos al integrarlo en herramientas automatizadas.
    \item \textbf{Rendimiento}: la lectura del problema es delegada a Unified Planning; para cargas repetidas grandes podría convenir cachear el resultado de \texttt{read\_problem} o parsear selectivamente los fluentes necesarios.
\end{itemize}

\subsection{Mejoras sugeridas}
\begin{itemize}
    \item Soportar identificadores de casilla genéricos usando expresiones regulares en \texttt{\_extraer\_coordenadas} para admitir índices multi-dígito.
    \item Añadir funciones utilitarias que extraigan la lista de casillas y relaciones de adyacencia desde \texttt{self.problem} para evitar recalcular adyacencias desde las acciones a lo largo de la búsqueda.
    \item Incorporar pruebas unitarias que verifiquen la correcta conversión de fluentes y el cálculo de la \textit{pagoda} para distintas variantes y tamaños de tablero.
\end{itemize}

\section{Descripción de los ficheros en el directorio PDDL}

En el directorio \texttt{PDDL/} se almacenan el dominio común del juego y varios problemas (configuraciones de tablero). A continuación se describe el contenido y las decisiones de diseño principales de cada fichero.

\subsection{\texttt{dominio\_senku.pddl}}
Fichero que define el dominio PDDL del juego Senku. Contiene:
\begin{itemize}
    \item Predicados: \texttt{(pieza ?c)} para casillas ocupadas, \texttt{(vacio ?c)} para casillas libres y \texttt{(adyacente ?c1 ?c2 ?c3)} para tripletas alineadas que representan un salto válido.
    \item Requerimientos: usa \texttt{:strips}, por tanto las acciones se expresan con precondiciones y efectos clásicos.
    \item Acción \texttt{saltar}: parámetros \texttt{?origen ?intermedia ?destino}, precondiciones que exigen la pieza origen, la pieza intermedia y destino vacío, además de la relación \texttt{adyacente}. Los efectos actualizan los predicados \texttt{pieza} y \texttt{vacio} apropiadamente.
\end{itemize}

Diseño notable: la relación \texttt{adyacente} se modela como predicado estático que encadena tres casillas formando una línea de salto; esto evita cálculos geométricos dentro del dominio PDDL pero genera muchos fluentes estáticos en los problemas.

\subsection{\texttt{problema\_senku\_prueba.pddl}}
Un problema de prueba con un tablero reducido y sólo tres piezas en posiciones concretas. Contiene:
\begin{itemize}
    \item Declaración de objetos: lista de identificadores de casilla con convención \texttt{cXY}.
    \item \texttt{(:init)}: todas las relaciones \texttt{adyacente} necesarias (horizontales y verticales) y un estado inicial reducido donde sólo unas pocas casillas contienen \texttt{pieza} y el resto son \texttt{vacio} según la configuración de prueba.
    \item \texttt{(:goal)}: estado objetivo específico (se indica una disposición final con una pieza en el centro del tablero de prueba).
\end{itemize}

Uso: sirve para depurar la lógica del dominio y los algoritmos de búsqueda en un caso pequeño y controlado.

\subsection{\texttt{problema\_senku\_tablero\_clasico.pddl}}
Problema que modela el tablero clásico de Senku. Estructura similar a la del fichero de prueba pero con el estado inicial completo del tablero clásico (todas las casillas ocupadas salvo el centro) y la definición exhaustiva de \texttt{adyacente} para el tamaño clásico.

Peculiaridad: el estado inicial coloca \texttt{vacio c44} para representar el hueco central; el objetivo exige dejar una única pieza en la posición central (`c44`).

\subsection{\texttt{problema\_senku\_tablero\_variante1.pddl}}
Variante del tablero con una geometría ligeramente distinta (añade filas/columnas adicionales y casillas intermedias como \texttt{c22}, \texttt{c62}, \texttt{c66}, ...). Contiene su propia lista de objetos, relaciones \texttt{adyacente} adaptadas y un estado inicial/objetivo acordes a la forma de la variante.

Propósito: evaluar el comportamiento del motor de búsqueda frente a tableros no estándar y comprobar la generalidad del modelado basado en \texttt{adyacente} y la convención \texttt{cXY}.

\subsection{\texttt{problema\_senku\_tablero\_variante3.pddl}}
Variante más grande que extiende filas hasta índices terminados en \texttt{3..8} introduciendo casillas como \texttt{c83}, \texttt{c84}, etc. Mantiene la misma estructura: objetos, muchas relaciones \texttt{adyacente}, estado inicial con casi todas las casillas ocupadas salvo un hueco central (\texttt{vacio c44}) y objetivo consistente con el centro ocupado.

Uso experimental: probar escalabilidad y comprobar que el código acepta dimensiones mayores sin cambios en el dominio.

\subsection{\texttt{problema\_senku\_tablero\_variante5.pddl}}
Variante irregular y más amplia que utiliza identificadores hasta dos dígitos en la segunda componente (por ejemplo \texttt{c95}) y una topología no rectangular regular. Incluye muchas relaciones \texttt{adyacente} y define un estado inicial y objetivo particular (en este caso, el objetivo mantiene una pieza en una casilla central distinta, e.g., \texttt{pieza c55}).

Notas sobre las variantes:
\begin{itemize}
    \item Todas las instancias comparten el mismo \texttt{dominio\_senku.pddl}, lo que facilita reusar el motor de planificación.
    \item La convención de nombres \texttt{cXY} es consistente en los ficheros, pero algunas variantes usan índices con más de una cifra por coordenada; esto obliga a que cualquier código que extraiga coordenadas (p. ej. \texttt{\_extraer\_coordenadas} en \texttt{src/pddl\_reader.py}) sea robusto frente a identificadores de longitud variable.
    \item El uso intensivo de \texttt{adyacente} genera muchos fluentes estáticos; en soluciones de búsqueda se recomienda filtrar estos fluentes para acelerar el cálculo de hashes y de movimientos aplicables.
\end{itemize}

\section{Descripción del módulo src/beam\_search.py}

\subsection{Propósito}
El módulo \texttt{src/beam\_search.py} implementa la variante de Búsqueda en Haz (Beam Search) adaptada al problema Senku. Proporciona la función principal \texttt{busqueda\_haz(problema, tamano\_haz, max\_profundidad)} que, dada una instancia de problema cargada por \texttt{unified\_planning}, explora los estados por niveles manteniendo únicamente los mejores \texttt{W} sucesores según una heurística ligera.

\subsection{Componentes principales}
\begin{itemize}
    \item \textbf{Heurística:} \texttt{heuristica\_senku(estado, cached\_keys)} utiliza la \textit{función pagoda} definida en \texttt{src/pddl\_reader.py} para asignar un valor a un estado (suma de pesos sobre casillas ocupadas). La implementación actual crea un objeto \texttt{SenkuPDDLReader} y llama a \texttt{funcion\_pagoda} con \texttt{variante=1}.
    \item \textbf{Función principal:} \texttt{busqueda\_haz} realiza las siguientes tareas:
        \begin{itemize}
            \item Filtra los fluentes relevantes y extrae las adyacencias usando utilidades de \texttt{src/basic\_search}.
            \item Usa \texttt{SequentialSimulator} de \texttt{unified\_planning} para generar estados sucesores aplicando la acción \texttt{saltar}.
            \item Mantiene un conjunto \texttt{visitados} con hashes (\texttt{hash\_estado}) para evitar ciclos.
            \item En cada nivel genera todos los sucesores, los ordena por heurística y conserva los primeros \texttt{tamano\_haz} para formar la frontera del siguiente nivel.
        \end{itemize}
    \item \textbf{Integración:} El módulo depende de \texttt{Nodo}, \texttt{_generar\_movimientos}, \texttt{_extraer\_adyacencias}, \texttt{_filtrar\_fluentes\_relevantes} y \texttt{hash\_estado} definidos en \texttt{src/basic\_search.py}.
\end{itemize}

\subsection{Decisiones de diseño}
\begin{itemize}
    \item \textbf{Simulación fiel a UP:} el uso de \texttt{SequentialSimulator} garantiza que la aplicación de acciones respeta la semántica del problema PDDL y evita reimplementar la dinámica.
    \item \textbf{Heurística basada en pagoda:} se reutiliza la función pagoda como medida rápida de calidad de estado; se opta por variantes sencillas para mantener el coste heurístico bajo.
    \item \textbf{Selección por nivel (beam):} ordenar los sucesores por heurística y truncarlos a \texttt{W} es una forma estándar de limitar la expansión en nodos.
\end{itemize}

\subsection{Limitaciones observadas}
\begin{itemize}
    \item \textbf{Creación repetida de lector PDDL:} \texttt{heuristica\_senku} instancia un \texttt{SenkuPDDLReader} en cada llamada, lo que es innecesario si sólo se usa \texttt{funcion\_pagoda}; conviene reutilizar una instancia o extraer la función pagoda estática.
    \item \textbf{Heurística débil:} la pagoda implementada es muy simple y puede no distinguir bien estados cercanos al objetivo, provocando pérdidas de soluciones si \texttt{tamano\_haz} es pequeño.
    \item \textbf{Gestión de parámetros de acción:} el código detecta diferentes formatos de acciones (objetos con atributo \texttt{action} o tuplas), lo que añade fragilidad ante cambios en la API de generación de movimientos.
    \item \textbf{Ausencia de control de tiempo/límite de nodos:} no existe timeout ni límite en número total de nodos expandidos (salvo profundidad), lo que puede bloquear búsquedas largas.
\end{itemize}

\subsection{Mejoras sugeridas}
\begin{itemize}
    \item Reutilizar una sola instancia de \texttt{SenkuPDDLReader} o extraer \texttt{funcion\_pagoda} a una utilidad pura para evitar sobrecoste en la heurística.
    \item Permitir pasar una heurística por parámetro a \texttt{busqueda\_haz} (por ejemplo una función que reciba \texttt{estado} y \texttt{cached\_keys}).
    \item Añadir memoización de heurística y un mecanismo de timeout/contador máximo de nodos.
    \item Mejorar la selección de sucesores usando estructuras más eficientes (p. ej. heap parcial) cuando el número de sucesores es grande en comparación con \texttt{tamano\_haz}.
    \item Documentar y estandarizar el formato de las acciones generadas por \texttt{_generar\_movimientos} para evitar múltiples ramas de código al extraer \texttt{act} y \texttt{params}.
\end{itemize}

\subsection{Ejemplo de uso}
Un uso mínimo desde Python sería:
\begin{verbatim}
from unified_planning.io import PDDLReader
from src.beam_search import busqueda_haz

reader = PDDLReader()
problema = reader.parse_problem('PDDL/dominio_senku.pddl', 'PDDL/problema_senku_tablero_clasico.pddl')
plan = busqueda_haz(problema, tamano_haz=100, max_profundidad=50)
print(plan)
\end{verbatim}


\end{document}